\documentclass[12pt]{article}
\usepackage{geometry,fancyhdr,xr,hyperref,ifpdf,amsmath,rcs,shortcuts}
\usepackage{lastpage,longtable,color,paunits,amsmath,smallsec}
\geometry{letterpaper,headsep=5mm,head=10mm,hmargin={20mm,30mm},bottom=10mm,top=5mm}


\pagestyle{fancy} 
\lhead{Atsc. 405} 
\chead{Notes entrainment (2006/3/15)} 
\rhead{page~\thepage/\pageref{LastPage}} 
\lfoot{} 
\cfoot{} 
\rfoot{}


\ifpdf
    \usepackage[pdftex]{graphicx} 
    \usepackage{hyperref}
    \pdfcompresslevel=0
    \DeclareGraphicsExtensions{.pdf,.jpg,.mps,.png}
\else
    \usepackage{hyperref}
    \usepackage[dvips]{graphicx}
    \DeclareGraphicsRule{.eps.gz}{eps}{.eps.bb}{`gzip -d #1}
    \DeclareGraphicsExtensions{.eps,.eps.gz}
\fi


\begin{document}

\begin{center}
{\bf
Atsc. 405\\
Notes on entrainment: W\&H Section 6.3\\
 }
\end{center}


\section{A repeat of the Section 6.3 derivation}
\label{sec:repeat-derivation}

In section 6.3 W\&H discuss in some detail the effect of the entrainment/mixing
of environmental air on cloud temperature.  Here's my version of the
derivation of (6.18) on p.~221, which I think is simpler.

In W\&H notation, a cloud parcel with mass $m - dm$, temperature $T^\prime$ and
saturation mixing ratio $r_s(T^\prime, p)$ ascends through an unsaturated environment 
of temperature $T$ and total mixing ratio $w$.  We want to find equations that
will predict the potential temperature $\theta^\prime$  and the liquid and vapor
mixing ratios $r_v^\prime$ and $r_l^\prime$ of the cloud parcel as a function of
height or time (assuming that we know the vertical velocity at all times).

From our earlier work in week 9 we know that the entropy $\phi$ and moist static energy
$h_m$ are both conserved for adiabatic processes and \textit{mix linearly}.  By 
\textit{linear mixing} we mean that if we have two blobs of air with
energies $h_m^\prime$, total mass $m-dm$ and  $h_m$ and total mass $dm$, then the
total energy of the mixture is just the sum of the two individual energies.
\textit{Question:  what kinds of systems violate this assumption and are non-linearly mixing?}.
Mathematically, that means that:

\begin{equation}
  \label{eq:mixture}
  m\,h_{m\,mix} = (m - dm) h_m^\prime + dm\,h_m
\end{equation}
Where the unprimed variables are for the environment, and the primed variables are for the cloud parcel.

Defining the change in specific parcel energy $dh_m^\prime$ as 
(energy after - energy before)/(final mass m):

\begin{equation}
  \label{eq:change}
  dh^\prime_{m} = (h_{m\,mix} - h_m^\prime)
\end{equation}
We can rearrange (\ref{eq:mixture}) to get a version of the first law for the
change in parcel energy due to the mixing process:

\begin{equation}
  \label{eq:change2}
\Delta h_m^\prime = q dt = \frac{dm}{m} (h_m - h_m^\prime) = \frac{dm}{m} \Delta h_m
\end{equation}
where $dm/m$ is the fraction of environmental air entrained into the cloud,
$\Delta h_m$ is the energy difference between cloud and environment and we've extended
our definition of heating to include energy changes caused by the physical transport
of mass (\textit{Question: what was our former definition of heating?}).

We know that the cloud parcel has to obey the second law:

\begin{equation}
  \label{eq:2nd}
  d\phi^\prime \geq \frac{q dt}{T^\prime}
\end{equation}

If we neglect the irreversible effects of mixing on entropy (\textit{why?}), and use the definition of
the equivalent potential temperature we've got:

\begin{equation}
  \label{eq:2ndb}
  d\phi^\prime = c_p \frac{d\theta_e^\prime}{\theta_e^\prime} = \frac{q dt}{T^\prime} = 
          \frac{dm}{m} \frac{\Delta h_m}{T^\prime}
\end{equation}
Where we've replaced the integrated heating $q dt$ by the entrainment using (\ref{eq:change2}).

We can use the definiton of $h_m$ to write:

\begin{equation}
  \label{eq:definition}
  \Delta h_m = (h_m - h_m^\prime) = (c_p T + l_v\, r + g z) - (c_p T^\prime + l_v\, r_s(T^\prime,p) + g z)
\end{equation}
Putting (\ref{eq:definition}) into (\ref{eq:2ndb}) and expanding $\theta_e$ gives
W\&H (6.18):

\begin{equation}
  \label{eq:whcorr}
\frac{d\theta_e^\prime}{\theta_e} =   \frac{d\theta^\prime}{\theta} + \frac{l_v\,dr_s(T^\prime,p)}{T^\prime} = 
   \left [ \frac{T^\prime - T}{T^\prime} + \frac{l_v}{c_p T^\prime} (r_s(T^\prime,p) - r) \right ]
   \frac{dm}{m}
\end{equation}

Why do I like this derivation better than W\&H? 

\begin{itemize}
\item It's shorter (7 equations vs. 12)
\item It's easier to interpret.  W\&H make a mistake when they say on the top of p.~221
that the change in saturation mixing ratio $dr_s$ comes from cooling due to ascent.  
We know this because we
got the same term in (\ref{eq:whcorr}) for a mixture that occured at constant pressure.
The reason that $r_s$ is decreasing is because the mixture temperature is colder than
the original parcel temperature due to evaporation, i.e.:

\begin{eqnarray}
  \label{eq:mix}
  dr_s &\approx& \frac{\partial r_s}{\partial T} (T_{mix} - T^\prime)\nonumber\\
  \frac{d\theta^\prime}{\theta^\prime} &=& \frac{dT^\prime}{T^\prime} - 
       \frac{R_d}{c_p}\frac{dp^\prime}{p^\prime} \approx \frac{T_{mix} - T^\prime}{T^\prime}
\end{eqnarray}
\textit{Question: What is the difference between $dT^\prime$ and $\Delta T^\prime$?}


\end{itemize}


\section{Vertical evolution of cloud properties}
\label{sec:vert-evol-cloud}

To actually solve (\ref{eq:whcorr}) we need to say something about
$dm/m$.  The simplest assumption is to assume that the cloud is mixing
like an \textit{entraining plume} with \textit{entrainment rate}
$\lambda$:

\begin{equation}
  \label{eq:plume}
  \frac{1}{m} \frac{dm}{dt} = \lambda
\end{equation}
Equation (\ref{eq:plume}) says that the bigger the parcel gets, the
more it mixes, possibly because the larger surface area provides more
of an interface to mix across.  You should convince yourself that
the solution to (\ref{eq:plume}) for constant $\lambda$ is:

\begin{equation}
  \label{eq:sol}
  m(t) = m_0 \exp ( \lambda  t)
\end{equation}

Given the vertical velocity $u=dz/dt$ we
can use the chain rule to write (\ref{eq:plume}) as:


\begin{equation}
  \label{eq:plume2}
  \frac{1}{m} \frac{dm}{dz} \frac{dz}{dt} = \frac{1}{m} \frac{dm}{dz} u = \lambda
\end{equation}
I can then use (\ref{eq:2ndb}) to get an equation I could integrate to
find the $\theta_e$ of the cloud as a function of height:


\begin{equation}
  \label{eq:2ndc}
 \frac{c_p}{\theta_e^\prime} \frac{d\theta_e^\prime}{dz} = 
         \frac{\Delta h_m}{T^\prime} \frac{1}{m}\frac{dm}{dz} = 
   \frac{\Delta h_m}{T^\prime} \frac{\lambda}{u} =
  \frac{\Delta h_m}{T^\prime} \hat{\lambda}
\end{equation}
Where $\hat{\lambda} = \lambda/u$ is the entrainment rate
in units of $m^{-1}$.

You should do the following:

\begin{itemize}
\item Show that an entrainment rate of $\hat{\lambda}= 7 \times
  10^{-4}\ \un{m^{-1}}$ implies that a cloud parcel would double its
  mass during an ascent of 1 km.  (But as noted in class, this
is a pretyy poor approximation of what happens in a real cloud with
saturated updrafts and downdrafts).

\item Show that by the same reasoning as above the total mixing ratio $r_T^\prime$
for the cloud parcel obeys an equation of the form:

\begin{equation}
  \label{eq:water}
  \frac{dr_T^\prime}{dz} = (r_{env} - r_T^\prime) \hat{\lambda}
\end{equation}

Using matlab we can integrate (\ref{eq:2ndc}) and (\ref{eq:water})
given an initial parcel $T^\prime,\, r_T^\prime,\, p^\prime$ and 
knowledge of the environmental ($T,\,w$) and $\hat{\lambda}$ as a
function of height to find the vertical profile of the cloud
equivalent potential temperature and total water.

\end{itemize}

\end{document}
